\documentclass[10pt,a4paper]{amsart}
\usepackage[margin=19mm]{geometry}
\usepackage{amsmath,amssymb,mathtools}
\usepackage[scr=rsfs,cal=euler]{mathalfa}
\usepackage{xfrac}
\usepackage[unicode,hidelinks,bookmarks=false]{hyperref}
\newcommand{\ie}{i.e.\ }
\newcommand{\cf}{cf.\ }
\newcommand{\resp}{respectively\ }
\newcommand{\Subsection}{\subsection}
\providecommand{\nobreakdash}{\nobreak\hbox{-}}
\setlength{\emergencystretch}{2em}
\multlinegap0pt
\usepackage{bm}
\usepackage{bbm}
\usepackage{dsfont}
\usepackage{xspace}
\usepackage{cancel}
\usepackage{mathtools}
\usepackage{cleveref}
\usepackage[shortlabels]{enumitem}
\usepackage{amsthm}
\makeatletter
\@ifundefined{theorem}{\newtheorem{theorem}{Theorem}}{}
\@ifundefined{corollary}{\newtheorem{corollary}[theorem]{Corollary}}{}
\@ifundefined{lemma}{\newtheorem{lemma}[theorem]{Lemma}}{}
\@ifundefined{prop}{\newtheorem{prop}[theorem]{Proposition}}{}
\makeatother
\newcommand{\Uc}{\mathcal{U}}
\newcommand{\abrac}[1]{\left\langle#1\right\rangle}
\newcommand{\cone}{\mathrm{Cone}}
\newcommand{\ft}{\mathrm{FT}}
\newcommand{\hf}{\mathfrak{h}}
\newcommand{\id}{\mathrm{id}}
\newcommand{\paren}[1]{\left( #1 \right)}
\newcommand{\sbim}{\mathbb{S}\textnormal{Bim}}
\newcommand{\set}[1]{\left \{ #1 \right \}}
\newcommand{\xequals}[1]{\overset{#1}{=\joinrel=\joinrel=}}
\title[Serre Duality and the Whitehead Link]{Serre Duality and the Whitehead Link}
\author{Cailan Li}
\date{Source: arXiv:2509.22133v1 (CC BY 4.0)}
\begin{document}
\maketitle

\noindent\emph{Source and licence.} Serre Duality and the Whitehead Link. Version arXiv:2509.22133v1 (CC BY 4.0), distributed under the Creative Commons Attribution 4.0 International licence, as recorded in the arXiv licence metadata for this exact version and shown on the abstract page https://arxiv.org/abs/2509.22133v1. Full rights notice: licenses/arxiv-2509.22133-CC-BY-4.0.txt.\par\smallskip

\noindent\textit{Editorial scope.} Counted unit: the Theorem whose source label is \texttt{relsertheorem} in the article (source0.tex lines 1085--1088, source label \texttt{relsertheorem}), delivered together with the article's own complete original proof of it (source0.tex lines 1089--1113, one \texttt{proof} environment of 25 lines). No mathematics is written, repaired or re-derived: statement and proof are the article's own text line for line. Required context retained uncounted: pis0 (lines 953--959); spliteq (lines 993--996); preprop (lines 998--1005); semicor (lines 1042--1049); tenequiv (lines 1065--1068). Every definition, display and label of those lines is retained unchanged. Omitted from this package and not redistributed: 1--952, 960--992, 997--997, 1006--1041, 1050--1064, 1069--1084, 1114--1577. Nothing inside a retained excerpt is omitted or altered: every retained body is delivered exactly as the article's lines are written, so no omission receipt of any kind accompanies this package. Citations: the retained excerpts carry no citation command at all, so no bibliography entry of the article is transcribed and no reference is redistributed. Reference adaptation: none was needed and none was made; every reference inside the delivered unit resolves inside it, so no unresolved reference or citation is produced. Printed slips kept literal: no printed word, symbol, hypothesis or number of the counted statement or of its proof was altered. Layout and class: the article's own class is \texttt{article} with packages including amsfonts, amsmath, stackengine, caption, amssymb, xfrac, xcolor, amsthm, fullpage, dsfont, enumerate, mathtools, float, tikz-cd; the wrapper is the lane's pinned \texttt{amsart} editorial wrapper at 10pt on A4, which restates the theorem-like environments the retained text needs and adds the article's own macro definitions that the retained text needs (\texttt{\textbackslash Uc}, \texttt{\textbackslash abrac}, \texttt{\textbackslash cone}, \texttt{\textbackslash ft}, \texttt{\textbackslash hf}, \texttt{\textbackslash id}, \texttt{\textbackslash paren}, \texttt{\textbackslash sbim}, \texttt{\textbackslash set}, \texttt{\textbackslash xequals}), with extra wrapper packages (bm, bbm, dsfont, xspace, cancel, mathtools, cleveref, enumitem). The delivered numbering is the unit's own. Assets: no retained span contains a figure environment, a graphics-inclusion command, a PDF-page inclusion, a table or a TikZ picture; no image file is copied into this package, the package declares no asset dependency and no image file is redistributed with it. Licence: version arXiv:2509.22133v1 of this article is distributed under the Creative Commons Attribution 4.0 International licence, as recorded in the arXiv licence metadata for this exact version and shown on the abstract page https://arxiv.org/abs/2509.22133v1. A full rights notice is delivered at licenses/arxiv-2509.22133-CC-BY-4.0.txt. Characters: every retained excerpt is pure ASCII, so the delivered engine, XeLaTeX with its default Latin Modern text font, reports no missing character.
\begin{corollary}
\label{pis0}
\begin{enumerate}[(i)]
    \item  For $1\le k\le \lfloor\frac{m}{2}\rfloor$,  $\pi^+_s\paren{F^\bullet_{(st)^k}}\simeq \pi^+_s\paren{F^\bullet_{(st)^ks}}\simeq 0$ in $\sbim(\hf, \abrac{t} )$.
    \item $\pi^+_s\paren{F_s}\simeq 0$
\end{enumerate}
\end{corollary}

\begin{equation}
\label{spliteq}
    F_s^2\xrightarrow{\id_{F_s}\otimes_R \psi_s} R\to \mathrm{Tot}\paren{F_s(-1)\xrightarrow{\alpha_s} \boxed{F_s(1)}} \xrightarrow{+1}
\end{equation}

\begin{prop}
\label{preprop}
Let $\overbrace{F_sF_t\ldots}^{k }$ be the Rouquier complex where we alternate between $F_s$ and $F_t$ for $k$ terms.
\begin{enumerate}[(i)]
    \item $\pi_s^+(\overbrace{F_s\ldots F_t}^{a }F_s^2\overbrace{F_t\ldots}^{b })\simeq \pi_s^+(\overbrace{F_s\ldots F_s}^{c }F_t^2\overbrace{F_s\ldots}^{d })\simeq 0$ for $|a-b|\ge 2$ and $|c- d|\ge 2$ and $a+b\le 2m-3$
    \item $\pi_s^+(\overbrace{F_sF_t\ldots}^{k })\simeq 0$ for all $2\le k \le 2m-3$.
\end{enumerate}
\end{prop}

\begin{corollary}
\label{semicor}
    Let \, $\Uc_t^{\pm}$ be the full triangulated subcategories spanned by the Rouquier complexes $\set{F_w}_{e,t\neq w\in I_2(m)}$ and $\set{F_w^{-1}}_{e,t\neq w\in I_2(m)}$, respectively. Then we have semi-orthogonal decompositions
    \begin{align*}
       K^b( \sbim(\hf, I_2(m) )) & \simeq \paren{ \Uc_t^{-} \to K^b( \sbim(\hf, \abrac{t} ))  } \\
       & \simeq \paren{  K^b( \sbim(\hf, \abrac{t} ))\to \Uc_t^{+}   }
    \end{align*}
\end{corollary}

\begin{lemma}
\label{tenequiv}
   Tensoring on the left (or right) with $\ft_{\set{s,t}/t}$ (see \cref{fteq} below) gives an equivalence of categories $\Uc_t^{-}\to \Uc_t^{+}$ with inverse $\ft_{\set{s,t}/t}^{-1}$. 
\end{lemma}

\begin{theorem}
\label{relsertheorem}
    $\pi_s^-(X^\bullet)\simeq \pi_s^+( \ft_{\set{s,t}/t}\otimes_R X^\bullet )\simeq \pi_s^+(X^\bullet\otimes_R \ft_{\set{s,t}/t} )$ for all $X^\bullet\in K^b(\sbim(\hf, I_2(m)))$.
\end{theorem}

\begin{proof}
It suffices to prove the first equivalence as the second proceeds similarly. When $m$ is odd, we have that 
   \begin{equation}
   \label{fteq}
       \ft_{\set{s,t}/t}=F^\bullet_{(st)^m}F_t^{-2}\simeq\overbrace{F_s\ldots F_t}^{m-1}\overbrace{F_s\ldots F_s}^{m} F_t^{-1}\xequals{braid}  \overbrace{F_s\ldots F_s}^{m-2}F_t^2\overbrace{F_s\ldots F_s}^{m-2}  
   \end{equation} 
   As in the proof of \cref{preprop} using \cref{spliteq} we obtain the distinguished triangle
   \[ \overbrace{F_s\ldots F_s}^{m-2 }F_t^2\overbrace{F_s\ldots F_s}^{m-2 }\to \overbrace{F_s\ldots F_t}^{m -3}F_s^2\overbrace{F_t\ldots F_s}^{m -3}\to  \mathrm{Tot}\paren{\overbrace{F_sF_t\ldots}^{2m-3}(-1)\to\overbrace{F_sF_t\ldots}^{2m-3 }(1) } \xrightarrow{+1}\]
   Using \cref{preprop} $(ii)$ the last term is homotopy equivalent to 0 after applying $\pi_s^+$, so we obtain the following chain of homotopy equivalences
   \begin{equation}
   \label{pifteq}
    \pi_s^+\paren{  \ft_{\set{s,t}/t}}=  \pi_s^+\paren{  \overbrace{F_s\ldots F_s}^{m-2 }F_t^2\overbrace{F_s\ldots F_s}^{m-2 }}\simeq \pi_s^+\paren{  \overbrace{F_s\ldots F_t}^{m-3 }F_s^2\overbrace{F_t\ldots F_s}^{m -3}}\simeq \ldots \simeq \pi_s^+\paren{R}=R  
   \end{equation} 
   Note that it is crucial that the start and ends of the expressions above end in $s$ as  $\pi_s^+(F_s)\simeq 0$ by \cref{pis0}$(ii)$ but $\pi_s^+(F_t)\not\simeq 0$! When $m$ is even, \cref{fteq} has $F_s^2$ in the middle (for the case $m=2$, $\ft_{\set{s,t}/t}=F_s^2$) but the start and end still consists of $F_s$ so the same arguments above also show \cref{pifteq} in this case. \\

   Using \cref{semicor}, $X^\bullet$ fits into the distinguished triangle
   \[  \pi_s^-(X^\bullet)\xrightarrow{\eta} X^\bullet \to C^-(X^\bullet)\xrightarrow{+1}  \]
   where $C^-(X^\bullet)=\cone(\eta)\in \Uc^-_t$. Tensoring with $ \ft_{\set{s,t}/t}$ on the left gives us the distinguished triangle
   \[  \ft_{\set{s,t}/t}\otimes_R \pi_s^-(X^\bullet)\to  \ft_{\set{s,t}/t}\otimes_R X^\bullet \to  \ft_{\set{s,t}/t}\otimes_R C^-(X^\bullet)\xrightarrow{+1}  \]
  $\ft_{\set{s,t}/t}\otimes_R C^-(X^\bullet)\in \Uc_t^+$ by \cref{tenequiv} so the last term is homotopy equivalent to 0 after applying the triangulated functor $\pi_s^+$. Thus by bilinearity and \cref{pifteq} we see that
  \begin{equation}
  \label{relsereq}
  \pi_s^-(X^\bullet)\simeq \pi_s^+\paren{\ft_{\set{s,t}/t}}\otimes_R \pi_s^-(X^\bullet)= \pi_s^+\paren{\ft_{\set{s,t}/t}\otimes_R \pi_s^-(X^\bullet)}=   \pi_s^+\paren{\ft_{\set{s,t}/t}\otimes_R X^\bullet} 
  \end{equation}
\end{proof}

\end{document}
